% Measured visits against area, on log-log axes, so that a power law is a
% straight line and its exponent is the slope.
\begin{figure}[htbp]
  \centering
  \begin{tikzpicture}
    \begin{loglogaxis}[
      width=0.86\textwidth, height=6.6cm,
      xlabel={cells $n$}, ylabel={nodes visited},
      log basis x=2, log basis y=10,
      legend pos=north west, legend cell align=left,
      legend style={font=\scriptsize, draw=black!30},
      grid=both, grid style={black!12},
      tick label style={font=\scriptsize}, label style={font=\small},
    ]
      \addplot[mark=*, wallred, thick] coordinates
        {(256,228) (1024,903) (4096,3543) (16384,13996) (65536,54293)};
      \addlegendentry{Dijkstra \quad slope $1.00$}
      \addplot[mark=square*, black!55, thick] coordinates
        {(256,68) (1024,173) (4096,536) (16384,1477) (65536,4348)};
      \addlegendentry{length of the route $d$}
      \addplot[mark=triangle*, doorgreen, thick] coordinates
        {(256,53) (1024,115) (4096,293) (16384,703) (65536,1588)};
      \addlegendentry{region tree \quad slope $0.68$}
      \addplot[dashed, black!45, no marks, domain=256:65536, samples=2]
        {sqrt(x)};
      \addlegendentry{$\sqrt{n}$, for reference}
    \end{loglogaxis}
  \end{tikzpicture}
  \caption{Both axes logarithmic, so a power law is a straight line and the
  exponent is its slope. Dijkstra runs parallel to $n$; the region-tree search
  runs at slope $0.68$, below even the length of the route it returns, because
  one visit to a corridor leaf covers a whole straight run.}
  \label{fig:growth}
\end{figure}
